%% -------------------------------------
%%
%% Appendix page example file.
%%
%% ---------------------
%% Change log:
%% 	2026-03-15 - initial version.
%%
%% ---------------------
%% Authors:
%% 	Santiago Husain Cerruti - hola dot santiago at proton dot me
%%
%% ---------------------
%% Filename: back-matter-appendix-calculos-complementarios.tex
%% Part of FIUBA-thesis project.
%% Copyright (c) 2026 Santiago Husain Cerruti.
%% Licensed under the MIT License.
%% See LICENSE file for full text.
% -------------------------------------
\chapter{Cálculos complementarios}\label{appCalcMisc}

\section{Rotaciones elementales}
En el libro de \textcite{espana_book} se desarrollan las expresiones de rotaciones alrededor de los ejes de un sistema cartesiano, llamadas \emph{rotaciones elementales de Euler}. Si el ángulo a rotar es $\phi$, las expresiones de las matrices de rotación son las siguientes:

\begin{IEEEeqnarray}{rCl}
\mathbf{R}_x\left(\phi\right) & = &
\begin{bmatrix}
1 & 0 & 0 \\
0 & \cos \left(\phi\right) & \sin \left(\phi\right) \\
0 & -\sin \left(\phi\right) & \cos \left(\phi\right) \\
\end{bmatrix}
\end{IEEEeqnarray}

\begin{IEEEeqnarray}{rCl}
\mathbf{R}_y\left(\phi\right) & = &
\begin{bmatrix}
\cos \left(\phi\right) & 0 & -\sin \left(\phi\right) \\
0 & 1 & 0 \\
\sin \left(\phi\right) & 0 & \cos \left(\phi\right)\\
\end{bmatrix}
\end{IEEEeqnarray}

\begin{IEEEeqnarray}{rCl}
\mathbf{R}_z\left(\phi\right) & = &
	\begin{bmatrix}
	\cos \left(\phi\right) & \sin \left(\phi\right) & 0 \\
	-\sin \left(\phi\right) & \cos \left(\phi\right) & 0 \\
	0 & 0 & 1
	\end{bmatrix}
\end{IEEEeqnarray}

